\section{Test Case 22: Repeated Dynamic Refinement and Capacity Growth}

Test Case 22 extends the hybrid refinement path of Test Case 21 from one
accepted Akima event to three. It retains Maxwell rheology, Yarin evaporation,
stochastic Platen aerodynamic drag, nozzle insertion, and ordinary numerical
anchor beads. The initial jet contains 400 elements over 8 cm of a 16 cm
nozzle-to-collector domain. The timestep is $2.5\times10^{-8}$ s and the run
ends after 15,800 steps.

The dedicated validation runner sets developer-only initial and growth
reserves of 20 entries. This forces every accepted remesh beyond the current
capacity, whereas production runs retain the normal 100-entry
\texttt{incnpjet} policy. Each event must release the old topology,
evaporation, Platen, Coulomb, and Gaussian-history mappings before the host
allocations change. The new state and repacked Gaussian history are then bound
once. The host constructs the target mesh and applies conservation, while
Akima coefficient construction and spline interpolation execute on the GPU.

With NVFORTRAN 24.3, the CPU events occur at steps 14,301, 14,944, and 15,719.
The active element counts change from 413 to 455, from 455 to 498, and from
499 to 536. Capacity grows successively from 420 to 477, then to 520, and
finally to 558. The complete-force oracle reproduces these CPU event steps and
topologies. The native NVIDIA A30 obtains the same active counts and capacity
sequence; its third event occurs at step 15,701 because of the different
direct-Coulomb accumulation order.

The three-event contract is intentionally a same-compiler NVFORTRAN CPU/GPU
comparison. A different compiler can perturb the stochastic bending
trajectory enough to accept a different number of marginal remeshes within
the same final time. Test Case 21 remains the portable single-event GFortran
validation.

All 81 active anchors retain their positions, velocities, stresses, reference
radii, and evaporation radii at every remeshing event. Reference volume,
evaporated volume, mass, and charge are conserved to approximately
$10^{-16}$. The 80-row native CPU/GPU statistical comparison passes with a
relative tolerance of 2.5 percent, while the complete-force oracle agrees with
the CPU within $6\times10^{-7}$ relatively.

An OpenACC transfer audit records four complete state downloads: three
accepted refinement events and the final shutdown. It records exactly three
Gaussian-history uploads, three topology rebinds, and three evaporation-state
rebinds. Ordinary timesteps and unsuccessful threshold scans do not transfer
the complete bead state.

The commands
\texttt{tests/refinement/run\_test22.sh nvfortran},
\texttt{tests/refinement/run\_test22.sh openacc}, and
\texttt{tests/refinement/run\_test22.sh force-oracle} run the CPU, native GPU,
and diagnostic oracle validations, respectively.
